\Needspace{18\baselineskip}
\section{Semantic examples and regression proofs}

The examples are small abstract machines used to check the meaning and
edge cases of the definitions. They do not formalize the paper's
lambda-calculus translation, copying UAM, or program-equivalence examples.

\Needspace{17\baselineskip}
\subsection{A countdown machine}

\begin{definition}[Countdown]
The countdown machine has state set $\mathbb N$, transition
$s\to t$ if $s=t+1$, and final predicate $s=0$.
\end{definition}
\lean{Examples}{11}{14}{Examples.countdown}
\begin{proof}[Verification of the machine law]
If $s=0$, an outgoing transition would imply $0=t+1$, impossible for
a natural number $t$. Thus final states are normal as required by
the machine structure.
\end{proof}

\begin{lemma}[Countdown is deterministic]
If $s\to a$ and $s\to b$, then $a=b$.
\end{lemma}
\lean{Examples}{16}{19}{Examples.countdown_deterministic}
\begin{proof}
The premises give $s=a+1$ and $s=b+1$. Cancellation of successor
gives $a=b$.
\end{proof}

\begin{lemma}[Canonical countdown execution]
For every $n\in\mathbb N$, the countdown machine has
$n\steps{n}0$.
\end{lemma}
\lean{Examples}{21}{24}{Examples.countdown_run}
\begin{proof}
Induct on $n$. For $n=0$, use the zero-step constructor. For $n+1$,
take the transition $n+1\to n$ and prepend it to the induction
hypothesis $n\steps{n}0$.
\end{proof}

\begin{proposition}[Countdown characterizes the cost relation]
For every cost preorder $Q$ and states $a,b\in\mathbb N$,
\[
 a\preceq_Q b\quad\Longleftrightarrow\quad a\,Q\,b.
\]
\end{proposition}
\lean{Examples}{27}{36}{Examples.countdown_refines_iff}
\begin{proof}
Suppose $a\preceq_Q b$. Apply it to the canonical successful execution
$a\steps{a}0$. It gives a successful execution of some length $m$
from $b$ and $a\,Q\,m$. Compare this execution with
$b\steps{b}0$. Determinism and final-execution uniqueness give
$m=b$, so $a\,Q\,b$.

Conversely assume $a\,Q\,b$. For any successful execution of length
$n$ from $a$, compare it with $a\steps{a}0$. Uniqueness gives
$n=a$. Choose $b\steps{b}0$ as the target execution; its cost
condition is $n\,Q\,b$, obtained by replacing $n$ with $a$.
\end{proof}

The file then proves three anonymous examples, each by the proposition
and the indicated numerical fact:
\begin{enumerate}
\item $5\preceq_I3$, because $3\leq5$.
\item $3\not\preceq_I5$, because $5\nleq3$.
\item $5\not\preceq_E3$, because $5\ne3$.
\end{enumerate}
\lean{Examples}{38}{39}{Examples: improvement from 5 to 3}
\lean{Examples}{41}{44}{Examples: no improvement from 3 to 5}
\lean{Examples}{46}{49}{Examples: unequal exact counts}
These verify that improvement uses the target count as the smaller
quantity. Reversing the intended cost orientation would invalidate the
first two examples.

\Needspace{17\baselineskip}
\subsection{A crash and an infinite loop}

\begin{definition}[Fault states and the fault machine]
There are four distinct states:
\[
 \mathsf{done},\quad\mathsf{crash},\quad\mathsf{stuck},\quad\mathsf{loop}.
\]
The only transitions are
\[
 \mathsf{crash}\to\mathsf{stuck},\qquad
 \mathsf{loop}\to\mathsf{loop}.
\]
Only $\mathsf{done}$ is final.
\end{definition}
\lean{Examples}{51}{53}{Examples.FaultState}
\lean{Examples}{55}{60}{Examples.fault}
The four names are constructors of a finite datatype with derived
decidable equality. Here the constructor $\mathsf{stuck}$ names a
particular state; $\stuck(s)$ remains the general predicate defined
earlier.
\begin{proof}[Verification of the machine law]
No transition has source $\mathsf{done}$, because both possible source
constructors are distinct from it. Consequently a final state has no
outgoing transition.
\end{proof}

\begin{lemma}[The fault machine is deterministic]
Any two transitions with the same source have the same target.
\end{lemma}
\lean{Examples}{62}{65}{Examples.fault_deterministic}
\begin{proof}
Case-analyze the two transition proofs. If both use the crash rule,
both targets are $\mathsf{stuck}$. If both use the loop rule, both
targets are $\mathsf{loop}$. Either mixed case would equate
$\mathsf{crash}$ with $\mathsf{loop}$ and is impossible.
\end{proof}

\begin{definition}[The crash--loop candidate]
Let $R(s,t)$ mean $s=\mathsf{crash}$ and $t=\mathsf{loop}$.
\end{definition}
\lean{Examples}{67}{67}{Examples.crashLoop}

\begin{lemma}[A simulation using the stuck alternative]
The crash--loop candidate is an $(E,E,E)$-counting simulation.
\end{lemma}
\lean{Examples}{71}{88}{Examples.crash_loop_simulation}
\begin{proof}
For its only related pair, the source is $\mathsf{crash}$, which is
not final; hence the finality condition is vacuous. Its only successor
is $\mathsf{stuck}$. Use the zero-step execution from that successor
to itself in the stuck alternative of simulation. It is normal because
neither transition rule starts there, and it is not final because it
differs from $\mathsf{done}$.
\end{proof}

The next anonymous example concludes
$\mathsf{crash}\preceq_E\mathsf{loop}$.
\lean{Examples}{90}{91}{Examples: crash refines loop}
\begin{proof}
Apply counting-simulation soundness to the preceding simulation, fault
determinism, the reasonable triple $(E,E,E)$, and the only pair of $R$.
\end{proof}
There is no requirement to match a successful source execution, because
none exists. Thus the result is consistent even with exact costs. The
example exercises the stuck branch of the implemented proof; it does
not prove that every conceivable up-to simulation proof of this same
refinement must use that branch.

\Needspace{23\baselineskip}
\Needspace{17\baselineskip}
\subsection{The zero-step converse case}

\begin{definition}[Visit specimens and endpoint maps]
The type of specimens has three distinct constructors
$\mathsf{entering},\mathsf{settled},\mathsf{finished}$, with derived
decidable equality. The endpoint maps into the countdown machine are
\[
 \begin{aligned}
 \ell(\mathsf{entering})&=2,& r(\mathsf{entering})&=1,\\
 \ell(\mathsf{settled})&=1,& r(\mathsf{settled})&=1,\\
 \ell(\mathsf{finished})&=0,& r(\mathsf{finished})&=0.
 \end{aligned}
\]
\end{definition}
\lean{Examples}{93}{95}{Examples.Visit}
\lean{Examples}{97}{100}{Examples.visitLeft}
\lean{Examples}{102}{105}{Examples.visitRight}

\begin{definition}[Visit specimen system]
Every visit specimen has search mode. The entering predicate holds
exactly of the constructor $\mathsf{entering}$. Every state is declared
rooted. Together with the endpoint maps above, these data define a
specimen system over countdown.
\end{definition}
\lean{Examples}{107}{126}{Examples.visitSystem}
\begin{proof}[Verification of the interface laws]
Rootedness preservation is immediate since rootedness is the true
predicate. Final-state agreement holds because the endpoint pairs are
$(2,1)$, $(1,1)$, and $(0,0)$: only in the third pair is either endpoint
zero, and then both are.

For context-only traversal, consider the specimen. The constructor
$\mathsf{entering}$ cannot satisfy the settled premise: its mode is
search and its entering predicate holds. At the constructor
$\mathsf{settled}$, both endpoints are $1$. Choose
$\mathsf{finished}$ and the countdown transitions $1\to0$ on both
sides; its search mode is remote. At $\mathsf{finished}$, both
endpoints are zero and neither has an outgoing transition, contradicting
the availability premise. These cases exhaust the datatype.
\end{proof}
This is a deliberately abstract instance. For example, all its modes are
search, including that of the pair of final states. The specimen interface
does not enforce a concrete UAM convention for final focus modes.

\begin{lemma}[Visit input-safety]
The visit system is input-safe for $(H^{-1},E)$.
\end{lemma}
\lean{Examples}{129}{137}{Examples.visit_safe}
\begin{proof}
The entering premise forces the input specimen to be
$\mathsf{entering}$. Use safety~II with output specimen
$\mathsf{settled}$ and $k=l=0$. The left execution is the single
countdown transition $2\to1$; the right execution has zero steps at
$1$. The output specimen is searching without entering and is therefore
settled. The cost requirement is $1\,H^{-1}\,0$, or $0\leq1$.
\end{proof}

\begin{lemma}[Visit converse robustness]
The opposite visit system is robust for $(H,E,E)$.
\end{lemma}
\lean{Examples}{139}{142}{Examples.visit_converse_robust}
\begin{proof}
Its mode is always search. The rewrite-mode premise of robustness is
therefore impossible, so every required implication holds vacuously.
\end{proof}

\begin{proposition}[The converse zero-step regression]
In countdown, $1\preceq_H2$.
\end{proposition}
\lean{Examples}{146}{149}{Examples.converse_zero_step_case}
\begin{proof}
Use converse state sufficiency for the visit system. Countdown is
deterministic, $(H,E,E)$ is reasonable, and the preceding two lemmas
supply its safety and converse robustness premises. The entering
specimen is remote, both its endpoints are rooted, and it witnesses
$\Lft(2,1)$. The theorem consequently gives $1\preceq_H2$.

In the underlying converse simulation construction, safety~II provides
the pair $(1,1)$ after one original left step and zero original right
steps. The prescribed right step is $1\to0$. Context traversal at
$(1,1)$ supplies $(0,0)$ and adds one transition to the original left
execution, now $2\to1\to0$. Thus the source terminates in one step
and the target in two, satisfying $1\,H\,2$. This is precisely the
zero-step subcase proved in the converse argument.
\end{proof}
